% =============================================================================
% BIS 527: Introduction to Health Data Science
% HOMEWORK SOLUTION TEMPLATE -- write your answers in this file.
%
% HOW TO USE THIS FILE
%   1. Copy it and rename it:  hw1-yourlastname.tex
%   2. Fill in \yourname, \yournetid, \hwnumber just below.
%   3. Fill in the two required blocks: collaborators and AI declaration.
%   4. Write your answer under each \part. Delete parts you did not attempt --
%      but say so in Problem 7 rather than staying silent.
%   5. Compile, then submit to Canvas the files the assignment asks for --
%      always hw1-yourlastname.pdf and hw1-yourlastname.tex, plus
%      hw1-yourlastname.py when the assignment has a programming problem.
%      Homework 1 has one, in Problem 1, so all three are required.
%
% This file compiles as handed to you. If it stops compiling, the last thing
% you typed is the culprit -- LaTeX errors are almost always local.
%
% NOTHING IN THIS COURSE IS HANDWRITTEN. That includes figures and diagrams --
% draw them with tikz, or plot them in Python and \includegraphics the result.
% Scans and photographs of written work are not accepted.
%
% Never used LaTeX? Open a free Overleaf account, upload this file, press
% Recompile. That is the whole workflow.
%
% The problem stubs below match HOMEWORK 1. For later assignments, delete
% them and add \problem / \part lines matching that assignment.
% =============================================================================

\documentclass[11pt]{article}

\usepackage{fullpage}
\usepackage{mathptmx}
\usepackage{latexsym}
\usepackage{graphicx}
\usepackage{amsmath, amssymb, amsthm, amsfonts}
\usepackage{xcolor}
\usepackage{listings}
\usepackage{algorithm}
\usepackage{algpseudocode}
\usepackage{booktabs}
\usepackage{enumitem}
\usepackage{hyperref}

\hypersetup{colorlinks=true, linkcolor=blue!50!black,
            citecolor=blue!50!black, urlcolor=blue!50!black}

% =============================================================================
% FILL THESE IN
% =============================================================================
\newcommand{\yourname}{Your Name}
\newcommand{\yournetid}{abc12}
\newcommand{\hwnumber}{1}

% =============================================================================
% Machinery -- you should not need to edit below this line.
% =============================================================================
\newcommand{\handout}[5]{
   \noindent
   \begin{center}
   \framebox{
      \vbox{
    \hbox to 6.28in { {\bf BIS 527: Introduction to Health Data Science} \hfill #2 }
       \vspace{4mm}
       \hbox to 6.28in { {\Large \hfill #5  \hfill} }
       \vspace{2mm}
       \hbox to 6.28in { {\it #3 \hfill #4} }
      }
   }
   \end{center}
   \vspace*{4mm}
}
\newcommand{\hwheader}{%
  \handout{\hwnumber}{Fall 2026}{\yourname}{NetID: \yournetid}
          {Homework \hwnumber{} --- Solutions}
}

% --- Required declaration blocks --------------------------------------------
\newcommand{\collaborators}[1]{%
  \noindent\fbox{\parbox{\dimexpr\textwidth-2\fboxsep-2\fboxrule\relax}{%
    \small\textbf{Collaborators.} #1}}
  \vspace{2mm}
}
\newcommand{\aideclaration}[1]{%
  \noindent\fbox{\parbox{\dimexpr\textwidth-2\fboxsep-2\fboxrule\relax}{%
    \small\textbf{AI use declaration.} #1}}
  \vspace{4mm}
}

% --- Problem / part structure -----------------------------------------------
\newcommand{\problem}[2]{%
  \vspace{4mm}
  \noindent{\large\bfseries Problem #1.} {\large #2}
  \par\vspace{1mm}\hrule\vspace{2mm}
}
\newenvironment{parts}{\begin{enumerate}[label=\textbf{(\alph*)}, leftmargin=*]}
                      {\end{enumerate}}

\newcommand{\inductionparts}{%
  \par\smallskip
  \noindent\textbf{Base case.}\par\vspace{18mm}
  \noindent\textbf{Inductive hypothesis.}\par\vspace{18mm}
  \noindent\textbf{Inductive step.}\par\vspace{30mm}
}

% --- A box for "I got stuck here". Using it earns full credit. ---------------
\newcommand{\stuck}[1]{%
  \par\vspace{1mm}
  \noindent{\color{red!50!black}$\triangleright$ \textbf{Where I got stuck:} #1}
  \par\vspace{1mm}
}

\newtheorem{claim}{Claim}
\newtheorem{definition}{Definition}

\lstdefinestyle{pythonstyle}{
  language=Python,
  basicstyle=\ttfamily\small,
  keywordstyle=\color{blue!50!black}\bfseries,
  commentstyle=\color{gray}\itshape,
  stringstyle=\color{red!50!black},
  numbers=left, numberstyle=\tiny\color{gray},
  showstringspaces=false, breaklines=true,
  frame=single, framesep=3pt, tabsize=4
}
\lstset{style=pythonstyle}

\begin{document}

\hwheader

% =============================================================================
% BOTH BOXES BELOW ARE REQUIRED. Do not delete either one.
% =============================================================================

\collaborators{%
  Name everyone you discussed these problems with. Write ``None'' if you
  worked alone.
}

\aideclaration{%
  Name the \textbf{tool}. Say \textbf{what it did} (drafted prose? checked
  algebra? wrote code? found references?). Say \textbf{what you verified, and
  how}. Say \textbf{what you did not verify}.
  \par\smallskip
  Claiming more checking than you actually did is worse than admitting the gap.
  An undeclared use, or an inaccuracy traced to one, scores zero.
  \par\smallskip
  If you used no AI tools, write: ``No AI tools were used.''
}

% =============================================================================
\problem{1}{Implementation and measurement}

\emph{Code in the \texttt{.py} file. Tables, figures and answers here.}

\begin{parts}
  \item Implementations and test results.
    \begin{center}
    \begin{tabular}{llccc}
      \toprule
      test & $t$ & expected & \texttt{linear} & \texttt{binary} \\
      \midrule
      first position    & & & & \\
      last position     & & & & \\
      interior position & & & & \\
      absent            & & & & \\
      $n = 0$           & & & & \\
      $n = 1$           & & & & \\
      duplicates        & & & & \\
      \bottomrule
    \end{tabular}
    \end{center}

  \item Comparison counts at $n = 15$, and the two worst cases.
    \begin{center}
    \begin{tabular}{l*{8}{c}}
      \toprule
      target $t$ & 0 & 1 & 2 & 3 & 4 & 5 & 6 & 7 \\
      \midrule
      \texttt{linear} & & & & & & & & \\
      \texttt{binary} & & & & & & & & \\
      \bottomrule
    \end{tabular}
    \end{center}
    Worst case, \texttt{linear}: \qquad Worst case, \texttt{binary}:
    \begin{center}
    \begin{tabular}{lcccccc}
      \toprule
      $n$ & $1$ & $2$ & $3$ & $4$ & $7$ & $8$ \\
      \midrule
      worst case, \texttt{binary} & & & & & & \\
      \bottomrule
    \end{tabular}
    \end{center}

  \item The $n$ at which recursive \texttt{linear} first fails, the error, and
    the reason; whether recursive \texttt{binary} fails at $n = 2^{20}$; and the
    two loop versions.
    \begin{lstlisting}[numbers=none]
def linear_loop(A, lo, hi, t):
    # your loop here

def binary_loop(A, lo, hi, t):
    # your loop here
    \end{lstlisting}

  \item Timings, repetitions used and how you chose them, and the figure.
    \begin{center}
    \begin{tabular}{lccccc}
      \toprule
      $n$ & $10^3$ & $10^4$ & $10^5$ & $10^6$ & $10^7$ \\
      \midrule
      \texttt{linear\_loop} & & & & & \\
      \texttt{binary\_loop} & & & & & \\
      \bottomrule
    \end{tabular}
    \end{center}
    % \begin{center}
    %   \includegraphics[width=0.7\textwidth]{timing-plot.png}
    % \end{center}

  \item Fitted slopes, expected slopes, what $\Theta(\log n)$ looks like on
    these axes, and any anomaly with its probable cause.

    Slope, \texttt{linear}: \qquad Slope, \texttt{binary}:
\end{parts}

% =============================================================================
\problem{2}{Completeness: does the procedure halt?}

\begin{parts}
  \item A measure for \texttt{linear} and one for \texttt{binary}: the quantity,
    why it is a non-negative integer, and why it strictly decreases.

    \emph{Measure for \texttt{linear}:} \qquad
    \emph{Measure for \texttt{binary}:}

  \item The procedure \texttt{down}:
    \begin{enumerate}[label=\textbf{(\roman*)}, leftmargin=*]
      \item the $k \geq 0$ on which it halts, and the arguments for one $k$ on
        which it does not;
      \item the requirement in the definition of a measure that fails, and the
        call at which it first fails;
      \item the one-line change, and a measure for your version.
    \end{enumerate}

  \item The modified \texttt{binary} does not halt: array, target, sequence of
    calls, and where the measure stops decreasing.
    \begin{center}
    \begin{tabular}{ccc}
      \toprule
      call & \texttt{mid} & measure \\
      \midrule
      1 & & \\
      2 & & \\
      3 & & \\
      \bottomrule
    \end{tabular}
    \end{center}

  \item Boundary cases: length $0$ and length $1$, for both procedures.

  \item Two procedures:
    \begin{enumerate}[label=\textbf{(\roman*)}, leftmargin=*]
      \item halts on every input, not correct;
      \item correct whenever it returns, does not always halt.
    \end{enumerate}
\end{parts}

% =============================================================================
\problem{3}{Correctness: does it return the right answer?}

\begin{parts}
  \item The claim for \texttt{linear}, as one precise sentence covering both
    outcomes.

  \item Proof for \texttt{linear}, by induction on $d$.
    \inductionparts

  \item The claim for \texttt{binary}, and its proof by induction on $d$. The
    step has three branches; handle all three.
    \inductionparts
    \stuck{If one branch defeated you, state which and how far you got.}

  \item The step that fails without sortedness, and an unsorted array of length
    $5$ on which \texttt{binary} returns \textsc{Not-Found} although $t$ is
    present.
    \begin{center}
    \begin{tabular}{cccc}
      \toprule
      call & \texttt{lo}, \texttt{hi} & \texttt{mid} & branch taken \\
      \midrule
      1 & & & \\
      2 & & & \\
      3 & & & \\
      \bottomrule
    \end{tabular}
    \end{center}

  \item Duplicates:
    \begin{enumerate}[label=\textbf{(\roman*)}, leftmargin=*]
      \item what each procedure returns, and whether the claim is violated;
      \item \texttt{binary} modified to return the smallest such index;
      \item the claim it satisfies, and which branch of the proof changes.
    \end{enumerate}
\end{parts}

% =============================================================================
\problem{4}{Time and space: what does the procedure cost?}

\begin{parts}
  \item Recurrences with base cases, why the floor, the solutions, the check
    against Problem~1(b), the values at $n = 10{,}000$, and the best cases.

    $T_{\text{lin}}(n) = $ \qquad $T_{\text{bin}}(n) = $
    \begin{center}
    \begin{tabular}{lccccccc}
      \toprule
      $n$ & $1$ & $2$ & $3$ & $4$ & $7$ & $8$ & $15$ \\
      \midrule
      counted in Problem~1(b) & & & & & & & \\
      your formula & & & & & & & \\
      \bottomrule
    \end{tabular}
    \end{center}

  \item $\mathbb{E}[C] = (n+1)/2$ for \texttt{linear}, in $\Theta$ terms, and
    whether averaging changes the choice.

  \item Depths and extra space; the account of the Problem~1(c) failure; the
    space of the loop versions; and the comparison table.
    \begin{center}
    \begin{tabular}{lccccc}
      \toprule
      & worst & best & space, recursive & space, loop & needs sorted \\
      \midrule
      \texttt{linear} & & & & & \\
      \texttt{binary} & & & & & \\
      \bottomrule
    \end{tabular}
    \end{center}

  \item Sorted linked list: running time of each in $\Theta$ terms, the
    operation whose cost changed, and which procedure changes more.

  \item The registry:
    \begin{enumerate}[label=\textbf{(\roman*)}, leftmargin=*]
      \item the two totals, $q^\ast$, and its values;
      \item cost of one insertion, locating separated from making room;
      \item the property a storage scheme would need for $O(\log n)$ throughout.
    \end{enumerate}
    \begin{center}
    \begin{tabular}{lccc}
      \toprule
      $n$ & $10^3$ & $10^6$ & $10^9$ \\
      \midrule
      $q^\ast$ & & & \\
      \bottomrule
    \end{tabular}
    \end{center}
\end{parts}

\end{document}
